% =============================================================================
% Chapter 1: Introduction (3-4 pages)
% =============================================================================

\section{Introduction}
\label{sec:introduction}

\subsection{Motivation}
\label{sec:motivation}

Many graphs of interest exhibit community structure: their vertices
form groups that are more densely connected internally than to the
rest of the graph. Examples include friendship circles in social
networks, functional protein complexes in biological systems, and
clusters of behaviorally similar customers in financial networks.
Recovering this structure from the graph alone, without external
labels, is the problem of \emph{community
detection}~\cite{fortunato2010community}, with applications across
network science and, more recently, in anti-fraud
analytics~\cite{wang2024quantum}.

The dominant practical approach is to maximize Newman's
modularity~\cite{newman2006modularity} through local greedy search,
as in the Louvain~\cite{blondel2008louvain} and
Leiden~\cite{traag2019leiden} algorithms. These methods are fast and
deliver high-quality partitions on graphs of millions of vertices.
However, they are local heuristics: they do not in general find the
global modularity optimum, and they suffer from the \emph{resolution
limit}~\cite{fortunato2007resolution}, which prevents them from
detecting small communities embedded in larger graphs. These
limitations motivate the search for alternative approaches.

One alternative formulates modularity maximization as a Quadratic
Unconstrained Binary Optimization (QUBO)
problem~\cite{lucas2014ising}. The partition is then a vector of
binary variables, and modularity is a quadratic form over those
variables. The problem itself remains NP-hard, but the QUBO
formulation opens it to a different class of solvers: classical
local search~\cite{ushijima2017graph}, quantum
annealing~\cite{negre2020detecting}, and Coherent Ising
Machines~\cite{wang2024quantum}. Wang et
al.~\cite{wang2024quantum} recently demonstrated this last approach
on real banking data, but their method depends on specialized analog
hardware that is not widely available.

A different line of work, due to Schuetz et
al.~\cite{schuetz2022combinatorial}, replaces the dedicated hardware
with an unsupervised graph neural network. Their Physics-Inspired
GNN (PI-GNN) parameterizes the binary assignment vector through a
continuous relaxation and trains the GNN to minimize the relaxed
QUBO loss on a single GPU, without any labeled data. The QIGNN
framework~\cite{qignn2026} extends PI-GNN with iterative refinement:
the prediction of one training epoch is concatenated with the input
features of the next, allowing the network to revisit and improve
its earlier output. Both PI-GNN and QIGNN have so far been validated
on two-class (single-bit-per-vertex) combinatorial problems --- MaxCut and Maximum
Independent Set --- where each variable is a single binary
indicator.

The community detection problem reduces to a single-bit-per-vertex
formulation only in the special case of two communities. With $k > 2$ communities each vertex must
be assigned a one-hot vector of length $k$, the QUBO matrix grows
from $n \times n$ to $nk \times nk$, and the one-hot constraint
needs to be enforced through a penalty term in the loss. Whether
the QIGNN approach extends cleanly to this setting is the question
this thesis sets out to answer. The investigation produces both an
extension of the QIGNN framework to the multi-class case and, where
the multi-class formulation runs into difficulties, a hierarchical
algorithm that recovers competitive performance by recursively
applying the more reliable bipartition formulation.


\subsection{Problem statement}
\label{sec:problem_statement}

Let $G = (V, E)$ be an undirected graph with $|V| = n$ vertices and
$|E| = m$ edges. A \emph{community structure} on $G$ is a partition
of the vertex set into $k$ disjoint, non-empty subsets, encoded as
a function $C: V \to \{1, \ldots, k\}$ that assigns each vertex a
community label. The community detection problem asks for the
partition that best reflects the underlying structure of the graph,
where ``best'' is measured by a quality function on partitions.

Newman's modularity~\cite{newman2006modularity} is the quality
function adopted throughout this thesis. It is defined as
\begin{equation}
    M(C) = \frac{1}{2m} \sum_{v, w \in V}
    \left( A_{vw} - \frac{k_v k_w}{2m} \right)
    \delta(C_v, C_w),
    \label{eq:modularity_intro}
\end{equation}
where $A_{vw}$ is the adjacency matrix entry between vertices $v$
and $w$, $k_v = \sum_w A_{vw}$ is the degree of vertex $v$, and
$\delta$ is the Kronecker delta. The optimization problem is then
$C^* = \arg\max_{C} M(C)$.

Modularity maximization is NP-hard~\cite{brandes2008modularity}, so
exact solution is infeasible at any practical scale and approximate
methods are required. This thesis assumes that the number of
communities $k$ is given as part of the input. Determining $k$
automatically is its own line of research, and fixing it lets the
focus stay on the optimization itself.


\subsection{Approach overview}
\label{sec:approach_overview}

The approach combines a QUBO formulation of the modularity objective
with an unsupervised graph neural network as the solver. Given a
graph $G$ and a target number of communities $k$, the procedure has
three steps: (i) construct a QUBO matrix $\Q$ that encodes the
modularity objective; (ii) train a GNN whose output, after softmax
normalization, parameterizes a relaxed solution $\Pb(\theta) \in
[0,1]^{n \times k}$; (iii) recover a discrete partition by taking
$\arg\max$ along the community dimension of $\Pb$. The GNN is
trained from scratch on the input graph by minimizing the relaxed
QUBO loss $\Lcal(\theta) = \pb_{\mathrm{flat}}^T \, \Q \,
\pb_{\mathrm{flat}}$. No labeled data is used.

For the QUBO construction the formulations of Wang et
al.~\cite{wang2024quantum} are used. The bipartition formulation
($k=2$) encodes community membership through a single binary
variable per vertex and reduces to the quadratic form $\xb^T (-\B)
\xb$, where $\B$ is the modularity matrix. The multi-class
formulation ($k > 2$) introduces $nk$ binary variables under a
one-hot encoding constraint, which is enforced softly through a
quadratic penalty term. Both formulations are implemented in this
work.

For the neural solver the QIGNN framework~\cite{qignn2026} is used
as the starting point. QIGNN combines GraphSAGE convolutional
blocks with iterative refinement of the previous epoch's prediction.
The architecture is adapted to multi-class output by replacing the
sigmoid head with a softmax layer over $k$ classes. To make the
multi-class loss tractable on graphs of practical size, a
memory-efficient training procedure is developed that computes the
loss directly from the modularity matrix and the probability matrix
without ever materializing the full $nk \times nk$ QUBO matrix.

The multi-class loss has a wide local attractor at the trivial
all-in-one assignment, where every vertex is placed into the same
community. An orthogonality regularization adapted from
MinCutPool~\cite{bianchi2020spectral} is integrated to discourage
this collapse by penalizing overlap between community columns. Its
behavior is evaluated empirically on graphs with both balanced and
imbalanced community structures.

Training uses the Adam optimizer with a time-dependent annealing
penalty that pushes the relaxed solution toward the corners of the
hypercube as training progresses. A schematic of the architecture
and training loop is given in Section~\ref{sec:methodology}.


\subsection{Scientific novelty}
\label{sec:novelty}

The scientific novelty of this thesis is summarized in the
following points.

\begin{enumerate}
    \item \textbf{A hierarchical decomposition algorithm for
    multi-class community detection.} A new algorithm is proposed
    that recursively applies the stable bipartition QUBO
    formulation through QIGNN to recover community structure for
    arbitrary $k$, using modularity gain as the natural stopping
    criterion. To the best of the author's knowledge, this is the
    first work to combine QUBO-based neural solvers with
    hierarchical graph decomposition for community detection. The
    hierarchical method, in its automatic modularity-gain mode,
    matches Louvain in modularity within $0.02$ on every benchmark
    in the study (including graphs with up to forty-two
    ground-truth communities) and strictly exceeds Louvain in
    normalized mutual information on four of them. An optional
    target-$k$ mode that takes the ground-truth $k$ as input
    trades a modularity gap on imbalanced graphs for higher NMI,
    matching or exceeding Louvain on ten of thirteen benchmarks.
    The modularity-gain criterion automatically discovers the
    number of communities, recovering Louvain's choice of $k$
    on graphs where Louvain's heuristic is itself reliable.

    \item \textbf{Extension of the QIGNN framework to multi-class
    community detection.} The QIGNN
    framework~\cite{qignn2026} and its predecessor
    PI-GNN~\cite{schuetz2022combinatorial} have so far been
    validated only on two-class problems such as MaxCut and Maximum
    Independent Set. To the best of the author's knowledge, this
    thesis is the first to apply the QIGNN paradigm to community
    detection with arbitrary $k$. The extension involves a softmax
    output layer over $k$ classes and the linearization of one-hot
    encoded QUBO variables.

    \item \textbf{Empirical demonstration that the bipartition
    formulation of QIGNN matches or exceeds Louvain in agreement
    with ground-truth community structure.} On diverse binary
    community detection benchmarks the unsupervised QIGNN-based
    solver reaches normalized mutual information at or above
    Louvain's. This result has not, to the author's knowledge,
    been previously reported.

    \item \textbf{Memory-efficient training procedure for the
    multi-class QUBO loss.} The proposed implementation computes
    the QUBO loss directly from the modularity matrix and the
    probability matrix, without materializing the full QUBO
    matrix. This reduces memory complexity from quadratic to
    linear in the linearized problem size and makes training
    feasible on graphs that the original PI-GNN and QIGNN
    implementations cannot handle.

    \item \textbf{Analysis of orthogonality regularization in the
    QUBO setting.} The orthogonality regularization originates in
    the supervised MinCutPool method~\cite{bianchi2020spectral}
    (and is also used by DMoN~\cite{tsitsulin2020graph}); its
    behavior in the unsupervised QUBO setting has not previously
    been studied. A systematic empirical comparison across diverse
    graph structures is provided, identifying when the
    regularization helps and when its built-in balance prior
    actively harms performance on imbalanced graphs.
\end{enumerate}


\newpage